Les activités d'une association sont financées grâce à des donateurs.

On réalise une estimation du nombre de donateurs de la manière suivante.

Chaque année,
\begin{itemize}[left=5mm]
	\item 80\,\% des donateurs de l'année précédente renouvellent leur don;
	\item 300 nouveaux donateurs effectuent un don.
\end{itemize}

En 2025, l'association comptait \np{1000} donateurs.

Pour tout entier naturel $n$, on modélise par $u_{n}$ le nombre de donateurs pour l'année $2025+n$.

Ainsi, on a $u_{0}=1000$.

\subsection*{Partie A}

\begin{enumerate}
	\item \begin{enumerate}
		\item Déterminer le nombre de donateurs prévu en 2026.

		\item Justifier que pour tout entier naturel $n$, on a :\quad $u_{n+1}=0,8 u_{n}+300$.
	\end{enumerate}

	\item On donne $u_{5} \approx \np{1336}$. Que représente cette valeur dans le contexte de l'exercice ?

	\item On considère la suite d'instructions Python ci-dessous :

\begin{lstlisting}[style=python,xleftmargin=5cm,xrightmargin=5cm,numbers=none]
u=1000
i=0
while u<=1400:
	u=0.8*u + 300
	i=i+1
\end{lstlisting}

	À la fin de l'exécution de ce script, la variable i contient la valeur 8.\\
	Comment peut-on l'interpréter dans le cadre de l'exercice?
\end{enumerate}

\subsection*{Partie B}

\begin{minipage}{11cm}
	On considère la suite $\left(v_{n}\right)$ définie sur $\N$ par \quad $v_{n}=u_{n}-\np{1500}$.

\begin{enumerate}
	\item Démontrer que $\left(v_{n}\right)$ est une suite géométrique.

	Préciser sa raison et son premier terme.
\end{enumerate}
\end{minipage}
\hfill
\begin{tblr}{colspec={X[c,4cm]},hlines, vlines}
	Aide au calcul :
	\[\dfrac{1200}{0,8}=\np{1500}\]
	\[0,8 \times \np{1500}=\np{1200}\]
\end{tblr}

\begin{enumerate}[start=2]
	\item \begin{enumerate}
		\item Exprimer $v_{n}$ en fonction de $n$.

		\item En déduire, pour tout entier naturel $n$, que : \quad $u_{n}=\np{1500}-500 \times 0,8^{n}$.
	\end{enumerate}
\end{enumerate}
